% Nouveau EGA V (cote 25): archive cote 25, site page 87 (PDF page 88).
% typescript. Draft compared: none.
% AI-transcribed by Claude Opus 5.5 (claude-opus-5-5, adaptive thinking, effort high), Message Batches API, batch ega25b (msgbatch_01A7n5YuXubUPNK1zsepi8RG), 2026-11-13; not yet checked by eye.
\scanpage{25/87}
\mspage{81}

suivant qu'on suppose  dim X \hand{$\geqslant$} m , ou  dim X \hand{$\leqslant$} m , voire dim X \hand{$<$} m.
Par la suite, nous regardons la $C_{\hand{y}}\hand{$\subset$} P_{k(\hand{y})}$\hand{)} , correspondant au point
générique \hand{$y$} de  Grass$^{m+1}$ , et nous dispensons  de faire l'interprétation
des résultats obtenus en termes de  ``presque tous les points'' ...

Tout d'abord, nous avons déjà noté \xout{\illegible{1}}dans 5.3. (et rattrapage du
cas général dans nº 12)  que $C_{\hand{y}}$  coupe ``régulièrement'' $X_{\hand{y}}$ , de façon
plus pré\xout{\unsure{s}}cise et plus générale, pour tout F quasi-cohérent sur X,
la section $\phi_{\hand{y}}^{\hand{(m+1)}}$ du \inserted[typed]{Module loc libre}\xout{fibré vectoriel} de rang  m+1  sur  $X_{k(\hand{y})}$  dont le
schéma des zéros est \over{C}{$C$}$_{\hand{y}}$ ; est F-régulière. En vertu de 10.2. , cela
implique par exemple que le morphisme  $\widetilde{X}(\hand{C_y})\to X_{k(\hand{y})}$ identifie $\widetilde{X}(\hand{C_y})$ au préschéma
déduit de $X_{k(\hand{y})}$ par éclatement de $f^{-1}(C_{\hand{y}})= X\times_P C_{\hand{y}}$ ;  \xout{\illegible{2}} dans le cas
où  dim f(X) \hand{$\leqslant$} m ,  on  aura même  $f^{-1}(C_{\hand{y}}) = \emptyset$ , et  par suite
$\widetilde{X}(C_{\hand{y}}) \xrightarrow{\sim} X_{k(\hand{y})}$  est un \uline{isomorphisme} , (et en fait, la \hand{\unsure{)}}projection conique
restreinte  est alors définie sur X tout entier à priori). Il se pose
ensuite la question de la dimension des fibres de $p_{C_{\hand{y}}}$ : $\widetilde{X}(C_{\hand{y}}) \longrightarrow Q(C_{\hand{y}})$,
et de la platitude de ce morphisme. On trouve:
\uline{Proposition} 14.1.  \xout{a) Si dim X} Supposons X irréductible,\over{\unsure{e}}{\unsure{ou}}plus
généralement, que  pour toute composante irréductible $X_i$ de X, la fibre
de $X_i$ en le point $f(x_i)$  ($x_i$ = point générique de $X_i$) ait une dimension d
(indépendante de i), ce qui a lieu par exemple \inserted{(avec d=0)} si f:X$\to$P est quasi-fini\hand{,} \margin{\deleted{avec d=0}}
Alors :     a)  Si  dim f(X) \hand{$>$} m, alors la dimension  des fibres de
$p_C$  : $\widetilde{X}(C_{\hand{y}}) \longrightarrow \over{Q}{Q}(C_{\hand{y}})$  sont toutes égales à  dim X - m .   b)  Si dim X
$\leqslant$ m, \inserted{et si les fibres non vides des $X_i$ sur P sont \unsure{tts} de dimen d, alors} les fibres de $p_C$ \inserted{non vides} \inserted{$\colon X_{k(y)}\to Q(C_y)$} sont toutes de dimension \over{0}{d}, \deleted{i.e. $p_C$ est quasi-%
fini} (donc $p_C$ est fini \inserted{\unsure{resp. quasi-fini}} \deleted{si X est propre sur \unsure{k$(y)$} et} \inserted{$f\colon X\to P$ est fini \unsure{resp. quasi-fini}}].
\margin{\unsure{resp. quasi-fini}}

On \deleted{se ramène aussitôt, en décomposant X en ses composantes irréd.,}
\deleted{au cas où X est irréductible}. Dans le cas a), \xout{si il y avait une fibre}
\xout{correspondant à un} \xout{xxxxxxxxxxxx} on sait déjà (je l'espère !) que pour \uline{tout}  point \hand{$y$} de

\margin{87}
